%=================================================================
% PLANTILLA -- Respuesta a los revisores -- Segundo Corte
% LabDST RoboCoffee, ITESO.
%
% Como usarla:
%   1. Copien esta carpeta junto a la carpeta de su paper revisado
%      (las dos carpetas deben ser HERMANAS) y cambien "N" por el
%      numero de su equipo, en el nombre de la carpeta y del .tex.
%   2. Ajusten la ruta de \LabDSTpaper: relativa, SIN extension,
%      entre comillas y en UNA sola linea.
%   3. En el paper, marquen cada cambio con \linelabel{ln:...} (texto)
%      o \label{sec:/fig:/tab:/eq:...} (secciones, figuras, tablas,
%      ecuaciones).
%   4. Escriban un \LabDSTcomment por CADA clave del PDF de
%      observaciones (R1.1, R1.2, ..., R4.n), en el mismo orden.
%   5. Compilen primero el paper (pdflatex -> bibtex -> pdflatex ->
%      pdflatex) y despues esta carta DOS veces (pdflatex x2).
%      Si aparece "??", revisen la ruta o el nombre de la etiqueta.
%
% La carta se escribe en ingles, como el paper. El comentario del
% revisor se copia TEXTUAL desde el PDF de observaciones.
%
% \LabDSTcomment{Clave}{Estado}{Ubicacion}{Comentario}{Respuesta}
%   Estado:    A = Addressed
%              P = Partially addressed
%              N = Not addressed (justified)
%   Ubicacion: \rln{x} -> "line NN"     \rlns{x}{y} -> "lines NN--MM"
%              \Rln, \Rlns (inicio de oracion), Section~\ref{sec:...},
%              Fig.~\ref{fig:...}, Table~\ref{tab:...}, \eqref{eq:...}
%   \LabDSTquote{...} cita el texto NUEVO del paper.
%
% Borren los comentarios <...> y todo lo que este entre < >.
%=================================================================
\documentclass{LabDSTresponse}
\LabDSTpaper{"../Equipo 3 - cuarto corte/Equipo 3 - cuarto corte"}

\begin{document}
	\LabDSTtitle{Fourth Cut}{Team 3}{Vision-Based Orientation Estimation and Correction for the
		SOFIA Modular Kinetic Surface}
	
	%-----------------------------------------------------------------
	\LabDSTsection{Acknowledgements}
	
	The authors thank the reviewers for their
	detailed evaluation of the second-cut
	manuscript. Their observations contributed
	to improving the technical documentation,
	clarity, and consistency of the proposed
	SOFIA system.
	
	The following responses address the
	reviewers' comments individually, indicating
	the modifications introduced and the aspects
	that remain subject to further development.
	All references to sections, figures, tables,
	and lines correspond to the revised manuscript
	accompanying this letter.
	
	\textbf{Main changes.} Falta hacer esto.
	
	\LabDSTsummary
	
	%-----------------------------------------------------------------
	\LabDSTreviewer{1}{Formato, LaTeX y cumplimiento de la plantilla}
	
	\LabDSTcomment{R1.1}{}{}
	{}
	{}
	
	\LabDSTcomment{R1.2}{}{}
	{}
	{}
	
	\LabDSTcomment{R1.3}{}{}
	{}
	{}
	
	\LabDSTcomment{R1.4}{}{}
	{}
	{}
	
	
	\LabDSTcomment{R1.5}{}{}
	{}
	{}
	
	
	%-----------------------------------------------------------------
	\LabDSTreviewer{2}{Redacción científica e idioma}
	
	\LabDSTcomment{R2.1}{}{}
	{}
	{}
	
	\LabDSTcomment{R2.2}{}{}
	{}
	{}
	
	%-----------------------------------------------------------------
	\LabDSTreviewer{3}{Contenido técnico: Materials and Methods y Design and Implementation}
	
	\LabDSTcomment{R3.1}{}{}
	{}
	{}
	
	\LabDSTcomment{R3.2}{}{}
	{}
	{}
	
	\LabDSTcomment{R3.3}{}{}
	{}
	{}
	
	\LabDSTcomment{R3.4}{}{}
	{}
	{}
	
	%-----------------------------------------------------------------
	\LabDSTreviewer{4}{Metodología D-Lab, resultados y validación}
	
	\LabDSTcomment{R4.1}{}{}
	{}
	{}
	
	\LabDSTcomment{R4.2}{}{}
	{}
	{}
	
	\LabDSTcomment{R4.3}{}{}
	{}
	{}
	
	\LabDSTcomment{R4.4}{}{}
	{}
	{}
	
	
\end{document}
